\chapter{时间序列、采样与周期信号}

\section{本章导读：周期不是最高峰，而是一个需要检验的假设}

时间序列的困难在于，数据值和观测时刻同样重要。两组亮度值即使完全相同，只要时间戳不同，就可能对应完全不同的物理解释。对天文学来说，这一点尤其关键，因为观测经常被昼夜、天气、季节、轨道和巡天策略切割成不均匀的窗口。

所以，本章不会把周期搜索教成“扫一堆周期，取分数最高的那个”。一个候选周期只是一个假设，它必须同时接受采样窗口、噪声、相位折叠图、次优候选和物理常识的检查。很多错误周期之所以危险，正是因为它们在离散采样点上也能看起来很像真的。

读这一章时，可以把每一次周期分析想成一份小型证据报告：观测覆盖了多久，候选周期范围如何设定，最高分和次高分差多少，相位折叠是否连贯，是否存在一天或季节尺度的混叠风险。这样的记录会比一个孤立的“最佳周期”更接近科学结论。

\section{学习目标}

学完本章后，你应该能够：

\begin{itemize}
    \item 理解时间序列数据和普通表格数据的根本差别在于观测顺序与采样节奏；
    \item 说清楚采样、观测窗口、噪声和混叠为什么会影响周期恢复；
    \item 从正弦周期模型出发，写出最小的周期搜索思路；
    \item 理解相位折叠为什么能把跨越多个周期的观测压缩到一个周期内部；
    \item 用一个极小光变曲线示例完成“读入 $\rightarrow$ 候选周期搜索 $\rightarrow$ 相位折叠 $\rightarrow$ 物理解读”的完整流程。
\end{itemize}

进入本章前，建议你已经完成 Part 0 的函数、循环和简单批处理训练，因为周期搜索几乎总是“把同一段计算重复跑很多次”的问题。

\section{训练目标与贯穿微型项目}

本章的微型项目是一条教学光变曲线的周期搜索。正文负责解释时间戳、采样、混叠、相位折叠和最小正弦模型；训练材料则要求你把这些概念变成一份候选周期记录，而不是只输出一个“最佳周期”。

\begin{examplebox}[最小周期搜索项目结构]
\begin{lstlisting}[style=terminal]
period_search_demo/
  data/
    lightcurve_period_demo.csv
  notebooks/
    search_period.ipynb
  scripts/
    search_demo_period.py
  figures/
    phase_folded_lightcurve.png
  results/
    period_record.txt
\end{lstlisting}
\end{examplebox}

项目交付应包含：时间单位、观测覆盖、候选周期范围、比较指标、最佳候选与次优候选、相位折叠图和混叠风险说明。这会直接训练后续模型评估中的一个核心习惯：不要只保存最高分，也要保存为什么这个候选可信。

\section{为什么时间序列会比静态表格更难}

在星表分析中，一行记录通常可以看成一个对象的一组属性；打乱行顺序以后，数据本身的物理含义大多不会消失。时间序列则完全不同。一个测量值不仅有大小，还有发生的时刻。把顺序打乱、把时间间隔忽略，往往就等于把最关键的结构丢掉了。

因此，时间序列分析至少同时面对三类问题：

\begin{enumerate}
    \item \textbf{信号问题}：源本身是否真的在变化；
    \item \textbf{采样问题}：观测时刻是否足够覆盖这种变化；
    \item \textbf{噪声问题}：测量误差、天气、仪器漂移和窗口函数是否会伪造结构。
\end{enumerate}

这也是为什么“把亮度对时间画出来”只是起点，而不是终点。

\section{时间序列的最小数学表示}

最简单的时间序列可以写成一组离散观测：

\[
\{(t_i, y_i, \sigma_i)\}_{i=1}^N,
\]

其中：

\begin{itemize}
    \item $t_i$ 是第 $i$ 次观测时刻；
    \item $y_i$ 是该时刻的测量值，例如亮度、流量或径向速度；
    \item $\sigma_i$ 是该点的不确定度。
\end{itemize}

如果要讨论周期信号，一个最常见的第一近似是：

\[
y(t) = C + A\sin\left(\frac{2\pi t}{P} + \phi\right) + \epsilon(t),
\]

其中 $C$ 是平均水平，$A$ 是振幅，$P$ 是周期，$\phi$ 是相位，$\epsilon(t)$ 则表示噪声和未建模结构。

\begin{defbox}[周期信号的最小模型]
很多天文周期信号并不严格是正弦波，但正弦模型仍然是入门阶段最有价值的第一近似，因为它足够简单，又能把周期、振幅和相位这些关键概念清楚地区分开。
\end{defbox}

\section{采样与 cadence：为什么不是观测得越多就一定越好}

初学者常把“样本点多”直接等同于“信息充分”，但时间序列里更关键的是\textbf{观测时刻如何分布}。如果所有点都挤在很短的时间内，即使点数很多，你依然不一定能恢复长周期；如果不同周期只在同一相位上被观测，信息也会非常有限。

通常我们把观测节奏称为 cadence。它影响：

\begin{itemize}
    \item 周期覆盖是否完整；
    \item 相邻点之间是否足够分辨快速变化；
    \item 是否容易产生一天、一轨道或一季节尺度的窗口偏差；
    \item 后续周期搜索是否出现混叠峰。
\end{itemize}

\paragraph{均匀采样直觉和不均匀采样现实并不一样}

很多学生第一次学周期信号时，会自然想到 Nyquist 采样直觉：采样频率至少要高于信号频率的两倍。这个思想在均匀采样情形下非常重要，但在天文时域数据里，观测时刻往往并不均匀：

\begin{itemize}
    \item 夜间观测会被白天中断；
    \item 地基观测会被天气和季节打断；
    \item 某些目标只在特定时间窗可见。
\end{itemize}

于是，真实问题不再只是“平均采样间隔是多少”，而是“观测窗口和信号本身如何共同决定可恢复的周期结构”。教学上必须让学生意识到：\textbf{不均匀采样下，时间戳本身就是数据内容的一部分。}

\section{混叠与窗口函数：为什么会找到“看起来也对”的假周期}

设真实周期是 $P_{\mathrm{true}}$，但观测只发生在某些离散时间上，那么数据中的结构实际上是“真实信号”和“观测窗口”共同作用的结果。这样一来，周期搜索曲线里常常会出现多个峰，其中只有一个更接近真实周期，其余则是混叠候选。

从直觉上说，混叠出现的原因是：\textbf{某些错误周期也能把离散采样点勉强对齐成一个看似平滑的折叠曲线。}

这提醒我们，周期搜索绝不是“找一个最大峰然后机械接受”，而应该同时检查：

\begin{itemize}
    \item 该周期下的相位折叠曲线是否真的连贯；
    \item 邻近候选周期是否给出相似质量；
    \item 采样窗口是否天然偏向某些频率。
\end{itemize}

\paragraph{为什么一个“最高峰”也可能不够}

在教学和真实项目里，周期搜索结果常常不会只给出一个候选，而是给出一串高峰。此时最危险的误区，是把最高峰直接当成最终答案。更稳妥的做法通常是：

\begin{itemize}
    \item 同时保留前几名候选周期；
    \item 比较它们的相位折叠形状是否稳定；
    \item 检查它们是否只是某个日周期、月周期或季节窗口的别名峰。
\end{itemize}

也就是说，周期搜索不是“读出一个数字”，而是“提出若干候选，再逐步排除假峰”。

\section{相位折叠：把长时间轴压缩回一个周期}

如果一个信号的周期是 $P$，那么任意时刻 $t_i$ 都可以映射成相位：

\[
\phi_i = \left(\frac{t_i}{P}\right) - \left\lfloor \frac{t_i}{P} \right\rfloor,
\]

其中 $\phi_i \in [0,1)$。这一步叫做\textbf{相位折叠}。它的思想很简单：若信号真的具有周期 $P$，那么跨越多个周期的观测应当在相位空间里重新对齐。

相位折叠的意义在于：

\begin{itemize}
    \item 它把长时间跨度压缩到一个周期内部；
    \item 它让我们更容易比较不同周期假设下的曲线是否平滑；
    \item 它常常比只看原始时间轴更容易暴露错误周期。
\end{itemize}

\paragraph{相位分箱为什么有用}

当折叠后的点仍然很散时，常见做法是在相位区间里再做分箱平均。这样做不是为了“把噪声抹掉”，而是为了更容易比较不同候选周期下的整体形状。

若把相位区间分成若干箱，第 $k$ 个箱中的平均值可以写成：

\[
\bar{y}_k = \frac{1}{n_k}\sum_{i \in \mathrm{bin}\,k} y_i,
\]

其中 $n_k$ 是该箱中的点数。配套 notebook 里的相位分箱摘要，本质上就是在做这件事。

\begin{examplebox}[最小相位分箱统计]
\begin{lstlisting}[style=pythonstyle]
def phase_bin_summary(folded_points, n_bins=5):
    bins = [[] for _ in range(n_bins)]
    for phase, flux in folded_points:
        index = min(int(phase * n_bins), n_bins - 1)
        bins[index].append(flux)
    return [
        (bin_index / n_bins, sum(values) / len(values), len(values))
        for bin_index, values in enumerate(bins)
        if values
    ]
\end{lstlisting}
\end{examplebox}

\section{从正弦模型到最小周期搜索}

如果把上面的周期模型展开成

\[
y(t) = C + a\sin\omega t + b\cos\omega t + \epsilon(t),
\qquad \omega = \frac{2\pi}{P},
\]

那么对于\textbf{固定}的候选周期 $P$，未知量只剩下 $C$、$a$ 和 $b$。这意味着我们可以把周期搜索拆成两层：

\begin{enumerate}
    \item 外层扫描一组候选周期 $P$；
    \item 内层在每个 $P$ 上拟合最优的 $C,a,b$，并比较残差或 $\chi^2$。
\end{enumerate}

若观测不确定度为 $\sigma_i$，则一个自然的评价量是

\[
\chi^2(P) = \sum_{i=1}^{N}
\frac{\left[y_i - \left(C + a\sin\omega t_i + b\cos\omega t_i\right)\right]^2}{\sigma_i^2}.
\]

最优周期通常对应较小的 $\chi^2$。这其实已经很接近周期图方法的基本思想。Lomb--Scargle 周期图之所以经典，就是因为它专门针对不均匀采样情形，把“给定频率下正弦拟合好不好”这件事写成了更系统的频域统计量。

\subsection{固定周期后，求的是线性系数}

当候选周期 $P$ 固定后，未知量不再是周期本身，而是系数 $C$、$a$ 和 $b$。这时问题可以写成一个很小的线性代数系统：

\[
\begin{bmatrix}
1 & \sin \omega t_1 & \cos \omega t_1 \\
1 & \sin \omega t_2 & \cos \omega t_2 \\
\vdots & \vdots & \vdots \\
1 & \sin \omega t_N & \cos \omega t_N
\end{bmatrix}
\begin{bmatrix}
C \\ a \\ b
\end{bmatrix}
\approx
\begin{bmatrix}
y_1 \\ y_2 \\ \vdots \\ y_N
\end{bmatrix}.
\]

这也是为什么很多周期搜索实现会先扫周期，再在每个候选周期上解一个最小二乘问题。对教学阶段来说，真正重要的是看见这一步：\textbf{周期是外层搜索变量，振幅和偏移是内层线性参数。}

\subsection{从线性系数回到振幅和相位}

把周期模型写成 $a\sin\omega t + b\cos\omega t$ 并不是为了改变物理含义，而是为了让固定周期下的拟合变成线性问题。拟合结束后，仍然可以回到更熟悉的振幅--相位形式：

\[
A\sin(\omega t+\phi)
= A\cos\phi\,\sin\omega t + A\sin\phi\,\cos\omega t.
\]

因此，

\[
a = A\cos\phi,\qquad b = A\sin\phi,
\]

反过来就有

\[
A = \sqrt{a^2+b^2},
\qquad
\phi = \mathrm{atan2}(b,a).
\]

这说明一个好的周期搜索结果不只给出周期 $P$，还应该能解释该周期下的平均亮度、振幅和相位。否则，我们只是在读出一个候选数字，而没有真正建立信号模型。

\begin{examplebox}[最小相位折叠函数]
\begin{lstlisting}[style=pythonstyle]
from math import floor

def phase_fold(times, period):
    return [(t / period) - floor(t / period) for t in times]
\end{lstlisting}
\end{examplebox}

\subsection{把“扫描周期”改写成“扫描频率”更自然}

虽然教学里我们更常直接说“扫描周期”，但从算法角度看，很多周期搜索其实更自然地是在扫描频率

\[
f = \frac{1}{P},
\qquad
\omega = 2\pi f.
\]

这样做的好处是：

\begin{itemize}
    \item 频率网格通常更容易均匀采样；
    \item 很多周期图方法本来就是频域对象；
    \item 邻近候选之间的差别在频率空间里往往更直观。
\end{itemize}

因此，周期和频率只是两种等价视角；真正重要的是你是否清楚自己扫描的网格是什么、分辨率够不够，以及网格是否已经覆盖了想找的时间尺度。

\subsection{候选周期网格怎么选}

周期搜索开始前，应该先写清楚候选范围和网格间隔，而不是随手给一组数字。一个最小判断可以从三个量出发：

\begin{itemize}
    \item 总时间跨度 $T=t_{\max}-t_{\min}$ 决定你对低频结构的分辨能力；
    \item 最短相邻观测间隔和典型 cadence 限制你能否看见快速变化；
    \item 科学问题本身决定哪些周期范围有物理意义。
\end{itemize}

在频率空间里，常用的粗略直觉是相邻候选频率间隔应小于

\[
\Delta f \sim \frac{1}{T},
\]

实际扫描时还常会再做几倍 oversampling，以免错过峰值附近的结构。但更密的网格并不自动带来更可信的周期，它只是让候选更细；最终仍要靠折叠曲线、残差和采样窗口共同判断。

\begin{warnbox}[搜索范围本身也是假设]
如果你只在 $0.5$ 到 $2$ 天之间搜索周期，就不可能发现 $10$ 天周期。周期搜索结果必须和候选范围一起报告，否则读者无法判断“没找到”到底是物理结论，还是搜索范围太窄。
\end{warnbox}

\begin{protip}[本章 notebook 的算法是教材化简版]
配套 notebook 不直接依赖科学计算重库，而是用一个极小的数据集和手写的最小线性代数步骤，演示“扫描候选周期并比较拟合残差”的核心思想。真正的 Lomb--Scargle 实现可以留到后续扩展或使用 \texttt{astropy} 时再系统展开。
\end{protip}

\section{一个极小教学光变曲线}

本章数据文件 \nolinkurl{lightcurve_period_demo.csv} 记录了一条教学版变星光变曲线。每一行包含：

\begin{itemize}
    \item \texttt{sample\_id}：目标标识；
    \item \texttt{time\_day}：观测时刻；
    \item \texttt{flux\_norm}：归一化亮度；
    \item \texttt{flux\_error}：每个点的测量误差；
    \item \texttt{split}：数据角色；
    \item \texttt{object\_class}：教学标签。
\end{itemize}

这个数据集有意保留了几个重要特征：

\begin{itemize}
    \item 采样时刻并不完全均匀；
    \item 亮度变化接近正弦，但并非每个点都严格落在曲线上；
    \item 时间跨度覆盖多个周期，适合做相位折叠演示；
    \item 数据规模很小，便于学生直接观察每一个测量点。
\end{itemize}

\section{相位折叠后还应该看什么}

得到一个候选周期并做相位折叠之后，还不意味着工作已经结束。至少还应再问：

\begin{itemize}
    \item 折叠后的点是否形成一条相对平滑的主轨迹；
    \item 是否只是在少数相位区间里有数据，而其余相位仍然空缺；
    \item 若把周期稍微改动一点，结构是否立刻变坏；
    \item 是否需要把相位再分箱，观察平均趋势而不是只看离散散点。
\end{itemize}

这一层检查很重要，因为它把“数学候选”进一步转成“是否像真实周期信号”的判断。

\section{候选周期应该怎样记录}

一个更像科研记录的周期搜索结果，不应只保存最佳周期，而应至少记录前几个候选。每个候选可以包含：

\begin{itemize}
    \item 候选周期或频率；
    \item 拟合残差、$\chi^2$ 或周期图功率；
    \item 相位覆盖是否充分；
    \item 折叠曲线是否平滑；
    \item 是否接近已知窗口周期或混叠位置。
\end{itemize}

这种记录方式有一个直接好处：当后续加入更多数据、改变误差模型或改用 Lomb--Scargle 实现时，你可以回头比较候选排名是否稳定，而不是只记得“上次算出来一个周期”。

\section{一个最小周期搜索结果报告模板}

周期搜索的报告不能只写“最佳周期是 \(P\)”。更好的教学写法应同时交代时间覆盖、搜索网格、候选稳定性和折叠诊断。例如：

\begin{quote}
我们使用归一化光变曲线，在 \(0.5\) 到 \(2.0\,\mathrm{day}\) 的候选周期范围内扫描最小正弦模型，并用加权残差比较候选。最佳候选周期约为 \(1.25\,\mathrm{day}\)，相位折叠后点列呈现单峰近似正弦结构；若把周期略微偏移，折叠曲线的相干性下降。该结果仍受当前时间跨度、采样窗口、误差模型和候选周期范围限制，尚不能排除邻近混叠或更复杂的非正弦变化。
\end{quote}

这段报告有四个关键部件：第一，说明搜索范围，而不是把算法当作全宇宙搜索；第二，说明比较标准，例如残差、\(\chi^2\) 或周期图功率；第三，说明折叠图是否支持候选；第四，主动保留采样窗口和混叠风险。到了最终项目包，这些内容会分别进入 Data card、Evaluation card、Interpretation card 和 Reproducibility card。

\section{配套 notebook 做了什么}

配套 notebook 走的是一条完整而克制的教学主线：

\begin{itemize}
    \item 读取极小光变曲线并检查时间覆盖；
    \item 计算均值、振幅范围和典型测量误差；
    \item 在一个有限周期网格上做最小正弦拟合；
    \item 选出最佳候选周期并计算对应的相位；
    \item 如果安装了 \texttt{matplotlib}，画出原始光变和相位折叠后的曲线；
    \item 最后比较最佳周期与若干次优候选，解释为什么周期恢复不能只看单一数字。
\end{itemize}

\section{常见错误}

\begin{warnbox}[时间序列分析中的典型问题]
\begin{itemize}
    \item 只看观测点个数，不看时间分布；
    \item 把相位折叠当成绘图技巧，而不是周期假设检验；
    \item 在存在多个混叠候选时，只接受“最高峰”而不检查折叠质量；
    \item 忽略测量误差，默认每个点一样可靠；
    \item 在时间跨度太短时过度声称发现了长周期。
\end{itemize}
\end{warnbox}

\section{AI 助手如何帮助你}

在这一章，AI 助手适合帮助你：

\begin{itemize}
    \item 把周期搜索代码整理成更清晰的函数结构；
    \item 检查相位计算公式和单位是否一致；
    \item 帮你解释为什么某些候选周期只是混叠；
    \item 协助你把“原始时间轴图”和“相位折叠图”的物理差别写清楚。
\end{itemize}

但一个周期是否真的可信、某个峰是否只是窗口效应、以及数据是否支持更复杂模型，仍然需要你自己做科学判断。

\section{练习}

\begin{exercisebox}[基础练习]
把最佳周期上下各偏移一个小量，重新计算相位并比较折叠曲线的平滑程度。请用语言说明你为什么认为原周期或偏移周期更合理。
\end{exercisebox}

\begin{exercisebox}[进阶练习]
从
\[
y(t) = C + A\sin\left(\frac{2\pi t}{P} + \phi\right)
\]
出发，推导出
\[
y(t) = C + a\sin\omega t + b\cos\omega t
\]
的等价形式，并写出 $a,b$ 与 $A,\phi$ 的关系。
\end{exercisebox}

\begin{exercisebox}[开放练习]
假设你要处理一条真实 TESS 或地基变星光变曲线。请写一个“周期恢复前检查清单”，至少包含：时间单位、缺测情况、典型误差、候选周期范围、混叠检查和相位折叠验证。
\end{exercisebox}

\section{本章小结}

时间序列分析的核心，不只是“数据按时间排好”，而是信号、采样和噪声三者同时作用的结果。本章通过一个极小光变曲线示例，建立了从周期模型、候选周期扫描到相位折叠检验的基本链条。只要这条链条真正建立起来，你就能更稳健地进入变星、凌星、径向速度乃至更复杂时域天文学任务。
